\section{A Liouville theorem for the walk}\label{sec:liouville}

We show that a bounded harmonic map for the walk of
Theorem~\ref{thm:walk}, subject to continuity on bounded sets in the tracial
$2$-norm, is left multiplication by its value at $1$.  The probabilistic
part of the proof produces free Haar unitaries such that, on words with
pairwise distinct unsigned indices, a single map acts by one of two left
multipliers according to the sign of the first letter.
Theorem~\ref{thm:rigidity}, proved in Section~\ref{sec:rigidity},
forces those multipliers to coincide for bounded linear maps with the
same continuity in any finite von Neumann algebra with a faithful
normal tracial state.

\subsection{Continuity and tracial ultrapowers}

\begin{definition}\label{def:ball-continuity}
Let $(N,\rho)$ be a finite von Neumann algebra with a faithful normal
tracial state.  A linear map $R:N\to N$ is \emph{bounded-ball
$2$-continuous} if
\begin{equation}\label{eq:ball-sequential}
 \sup_j\norm{x_j}<\infty,\qquad \norm{x_j}_2\longrightarrow0
 \quad\Longrightarrow\quad
 \norm{R(x_j)}_2\longrightarrow0.
\end{equation}
Here and below $\norm{x}_2=\rho(x^*x)^{1/2}$ when the ambient algebra is
$(N,\rho)$.
\end{definition}

The following lemma records the continuity needed after each of the two
ultrapowers in the argument.  It applies without a separability
assumption.

\begin{lemma}\label{lem:ultrapower}
Let $(N,\rho)$ be as in Definition~\ref{def:ball-continuity}, and let
$R:N\to N$ be bounded and linear.  Condition~\eqref{eq:ball-sequential}
is equivalent to the following uniform assertion: for every $r<\infty$
and $\varepsilon>0$ there is $\delta>0$ such that
\begin{equation}\label{eq:ball-modulus}
 \norm{x}\le r,\quad \norm{x}_2<\delta
 \quad\Longrightarrow\quad \norm{R(x)}_2<\varepsilon.
\end{equation}
If this condition holds and $\omega$ is a free ultrafilter on $\N$, then
\begin{equation}\label{eq:ultrapower-map}
 R^\omega:N^\omega\longrightarrow N^\omega,
 \qquad R^\omega([(x_j)]_\omega)=[(R(x_j))]_\omega,
\end{equation}
is a well-defined bounded linear map with
$\norm{R^\omega}\le\norm{R}$.  It is bounded-ball $2$-continuous for
the ultrapower trace.  In particular, for every free ultrafilter $\nu$,
the map $(R^\omega)^\nu$ on $(N^\omega)^\nu$ has the same properties.
\end{lemma}

\begin{proof}
If~\eqref{eq:ball-modulus} fails for some $r,\varepsilon$, choose
$x_j$ with $\norm{x_j}\le r$, $\norm{x_j}_2<1/j$, and
$\norm{R(x_j)}_2\ge\varepsilon$.  This contradicts
\eqref{eq:ball-sequential}; the converse is immediate.  Applying the
uniform assertion to differences gives uniform $2$-norm continuity on
each fixed operator norm ball.

Use the tracial ultrapower and its ideal $\mathcal I_\omega$ from
\eqref{eq:tracial-ultrapower}, now for $(N,\rho)$.  Its tracial norm is
\[
 \norm{[(x_j)]_\omega}_2=\lim_{j\to\omega}\norm{x_j}_2.
\]
By~\eqref{eq:ball-modulus}, the coordinate map induced by $R$ preserves
$\mathcal I_\omega$.  Applying this observation to the difference of two
bounded representatives proves that~\eqref{eq:ultrapower-map} is
well-defined.  The quotient norm gives
\[
 \norm{R^\omega(X)}
 \le \norm{R}\inf_{[(x_j)]_\omega=X}\sup_j\norm{x_j}
 =\norm{R}\norm{X}.
\]

To transfer the modulus, fix $r<\infty$ and $\varepsilon>0$, and choose
$\delta$ in~\eqref{eq:ball-modulus} for the radius $r+1$ and the output
tolerance $\varepsilon/2$.  If $\norm{X}\le r$, the quotient norm
allows a representative with $\sup_j\norm{x_j}\le r+1$.  When
$\norm{X}_2<\delta/2$, the set
\[
 \{j:\norm{x_j}_2<\delta\}
\]
belongs to $\omega$.  Consequently
\[
 \norm{R^\omega(X)}_2
 =\lim_{j\to\omega}\norm{R(x_j)}_2
 \le\varepsilon/2<\varepsilon.
\]
Thus $\delta/2$ is a valid input tolerance for $R^\omega$ on the
radius-$r$ ball.  The proof used no separability of $N$ and applies again
with $N^\omega$ in its place, proving the last assertion.
\end{proof}

\subsection{The harmonic map and its martingale}

The probabilistic step uses the classical passage from harmonic
functions to martingales along random products; see
\cite[p.~415]{Furstenberg1963}. Here the function is
$D(g)=T(g)g^*$ with values in $L^2(M,\tau)$, and the required
convergence follows directly from orthogonality of Hilbert-space
martingale increments.

\begin{theorem}\label{thm:liouville}
Let $(M,\tau)$ be a type $\mathrm{II}_1$ von Neumann algebra with
separable predual and a faithful normal tracial state, and let $\mu$ be
the measure supplied by Theorem~\ref{thm:walk}.  Suppose that
$T:M\to M$ is bounded, linear, and bounded-ball $2$-continuous, and that
\begin{equation}\label{eq:harmonic-map}
 T(x)=\E_{U\sim\mu}\bigl[T(xU)U^*\bigr]
 \qquad(x\in M).
\end{equation}
Then
\begin{equation}\label{eq:liouville-conclusion}
 T(x)=T(1)x\qquad(x\in M).
\end{equation}
\end{theorem}

\begin{proof}
Write $C=\norm{T}$ and $H=L^2(M,\tau)$.  Expectations in
\eqref{eq:harmonic-map} are norm-convergent sums, since $\mu$ is
countably supported and $\norm{T(xU)U^*}\le C\norm{x}$.  For
$g\in\U(M)$ put
\[
 D(g)=T(g)g^*.
\]
Let $U_1,U_2,\ldots$ be independent with law $\mu$, and set
$W_0=1$, $W_n=U_1\cdots U_n$, and
$\mathcal F_n=\sigma(U_1,\ldots,U_n)$.  The $H$-valued random
variables $Z_n=D(W_n)$ have countable range and satisfy
$\norm{Z_n}_2\le\norm{Z_n}\le C$.  Moreover,
\begin{align*}
 \E[Z_{n+1}\mid\mathcal F_n]
 &=\E\bigl[T(W_nU_{n+1})U_{n+1}^*W_n^*
                  \mid\mathcal F_n\bigr] \\
 &=T(W_n)W_n^*=Z_n.
\end{align*}
Thus $(Z_n)$ is a martingale in $L^2(\Omega;H)$.  Its increments are
orthogonal, and
\[
 \norm{Z_0}_2^2+
 \sum_{j=0}^{n-1}\E\norm{Z_{j+1}-Z_j}_2^2
 =\E\norm{Z_n}_2^2\le C^2.
\]
It follows that $Z_n$ converges in $L^2(\Omega;H)$ to a random variable
$Z_\infty$, and continuity of conditional expectation gives
$Z_n=\E[Z_\infty\mid\mathcal F_n]$.  In particular, with
\begin{equation}\label{eq:martingale-error}
 e_k=\E\norm{Z_\infty-Z_k}_2^2\longrightarrow0,
\end{equation}
orthogonality gives the uniform Cauchy estimate
\begin{equation}\label{eq:martingale-cauchy}
 \E\norm{Z_n-Z_k}_2^2=e_k-e_n\le e_k
 \qquad(n\ge k).
\end{equation}

The operator norm ball
$K_C=\{x\in M:\norm{x}\le C\}$ is closed in $H$.  Indeed, from a
$2$-norm-convergent sequence in $K_C$, ultraweak compactness provides
an ultraweakly convergent subnet with limit in $K_C$.  Testing against
$\tau(y^*\,\cdot)$, for $y\in M$, identifies this limit with the
$H$-limit, since $M$ is dense in $H$.  A subsequence of $Z_n$ converges
almost surely in $H$, so $Z_\infty\in K_C$ almost surely.

If $Z_\infty=c$ almost surely, conditional expectation gives
$Z_0=Z_1=c$ almost surely.  Hence $c=T(1)$ and $D(u)=c$ at every atom
$u$ of $\mu$.  These atoms contain a $2$-norm-dense subset of
$\U(M)$ by Theorem~\ref{thm:walk}.  For any $u\in\U(M)$, choose
atoms $u_j$ with $\norm{u_j-u}_2\to0$.  Bounded-ball continuity gives
$T(u_j)\to T(u)$ in $2$-norm, while
$\norm{cu_j-cu}_2\le\norm{c}\norm{u_j-u}_2\to0$.  Since
$T(u_j)=cu_j$, we obtain $T(u)=cu$.  Every element of a unital complex
$C^*$-algebra is a finite linear combination of unitaries, proving
\eqref{eq:liouville-conclusion} in this case.

Suppose, for a contradiction, that $Z_\infty$ is not deterministic.
Its law is a Borel probability measure on the separable Hilbert space
$H$, so its support contains two distinct points $a,b$.  Because
$K_C$ is closed and has full measure, $a,b\in K_C\subset M$.

\medskip\noindent
\emph{Forward and reverse endpoints.}
We claim that, for every $\delta>0$, there are numbers
$\eta_\delta>0$ and $N_\delta$ such that, whenever $n\ge N_\delta$
and $G$ has law $\mu^{*n}$,
\begin{equation}\label{eq:two-endpoints}
 \Pr\bigl(\norm{D(G)-a}_2<\delta,
          \ \norm{D(G^{-1})-b}_2<\delta\bigr)
 \ge\eta_\delta.
\end{equation}
Realize $G=U_1\cdots U_n$, let $h=\lfloor n/2\rfloor$, and put
\[
 A_n=D(U_1\cdots U_h),\qquad
 B_n=D(U_n^{-1}\cdots U_{n-h+1}^{-1}).
\]
The two variables use disjoint sets of increments and are independent.
Symmetry of $\mu$ implies that both have the law of $Z_h$.  Apply
\eqref{eq:martingale-cauchy} to the forward string and separately to
the reversed inverse string, which is also a string of independent
increments with law $\mu$.  It gives
\begin{equation}\label{eq:endpoint-replacement}
 \E\norm{D(G)-A_n}_2^2\le e_h,
 \qquad
 \E\norm{D(G^{-1})-B_n}_2^2\le e_h.
\end{equation}
For fixed $\delta$, the support property gives
\[
 p_a=\Pr(\norm{Z_\infty-a}_2<\delta/3)>0,
 \qquad
 p_b=\Pr(\norm{Z_\infty-b}_2<\delta/3)>0.
\]
Mean-square convergence implies, for all sufficiently large $h$,
\[
 \Pr(\norm{Z_h-a}_2<2\delta/3)\ge p_a/2,
 \qquad
 \Pr(\norm{Z_h-b}_2<2\delta/3)\ge p_b/2.
\]
Independence of $A_n,B_n$, the triangle inequality, and Markov's
inequality therefore bound the probability in
\eqref{eq:two-endpoints} from below by
\[
 \frac{p_ap_b}{4}-\frac{18e_h}{\delta^2}.
\]
For large $n$ this is at least $p_ap_b/8$, proving the claim.  The
independence used here is that of the disjoint prefixes $A_n,B_n$;
the variables $G$ and $G^{-1}$ remain dependent.

\medskip\noindent
\emph{Concatenation of independent blocks.}
Let $G_1,\ldots,G_m$ be independent with law $\mu^{*n}$, and consider
\[
 v(G)=G_{i_1}^{\sigma_1}\cdots G_{i_t}^{\sigma_t},
 \qquad \sigma_j\in\{1,-1\},
\]
where $i_1,\ldots,i_t$ are pairwise distinct.  Each inverse block can
be represented by reversing and inverting its own increments.  By
symmetry, concatenating these distinct signed blocks is a length-$tn$
walk with independent increments of law $\mu$.  In particular,
\[
 \bigl(D(v(G)),D(G_{i_1}^{\sigma_1})\bigr)
 \quad\hbox{has the law of}\quad (Z_{tn},Z_n).
\]
Consequently
\begin{equation}\label{eq:independent-blocks}
 \E\norm{D(v(G))-D(G_{i_1}^{\sigma_1})}_2^2\le e_n.
\end{equation}
The restriction is on unsigned indices: using both $G_i$ and $G_i^{-1}$
would reuse a block and would not justify this estimate.

\medskip\noindent
\emph{Finite requirements and diagonal selection.}
Fix an integer $m\ge1$ first, and use tolerance $1/m$ in
\eqref{eq:two-endpoints}.  For $m$ independent length-$n$ products,
let $E_{m,n}$ be the event
\[
 \norm{D(G_i)-a}_2<1/m,\qquad
 \norm{D(G_i^{-1})-b}_2<1/m
 \quad(1\le i\le m).
\]
Independence between the paths gives numbers $c_m>0$ and $N_m$ such
that
\begin{equation}\label{eq:conditioning-mass}
 \Pr(E_{m,n})\ge c_m\qquad(n\ge N_m).
\end{equation}
For example, one may take $c_m=(\eta_{1/m})^m$.

Let $\mathcal Q_m$ be the finite set of nonidentity reduced words of
length at most $m$ on the symbols $x_1,\ldots,x_m$.  Let
$\mathcal V_m$ be the finite set of all nonempty signed words on these
symbols with pairwise distinct unsigned indices.  Along the times
$n_\ell$ in Theorem~\ref{thm:walk}, the probability that any of the
requirements
\begin{align}
 |\tau(q(G_1,\ldots,G_m))|&<1/m
       &&(q\in\mathcal Q_m),\label{eq:selected-moments}\\
 \norm{D(v(G))-D(G_{i_1}^{\sigma_1})}_2&<1/m
       &&\bigl(v=x_{i_1}^{\sigma_1}\cdots x_{i_t}^{\sigma_t}
                   \in\mathcal V_m\bigr)\label{eq:selected-blocks}
\end{align}
fails tends to zero.  For the first line, this is a finite union of
the moment convergence statements in Theorem~\ref{thm:walk}.  For the
second line, the union bound and~\eqref{eq:independent-blocks} give
the upper bound $|\mathcal V_m|m^2e_{n_\ell}$.

Choose $\ell_m>\ell_{m-1}$, after fixing $m$, so large that
$n_{\ell_m}\ge N_m$ and this combined failure probability is less
than $c_m/2$.  Then the intersection of $E_{m,n_{\ell_m}}$ with all
the requirements~\eqref{eq:selected-moments}--\eqref{eq:selected-blocks}
has probability at least $c_m/2>0$.  Choose a deterministic tuple
\[
 (u_1^{(m)},\ldots,u_m^{(m)})\in\U(M)^m
\]
from that intersection.  The constants $c_m$ may tend to zero as fast
as they please: the choice of $\ell_m$ is made separately for each
fixed $m$, after the positive lower bound has been obtained.

Write $a_1=a$ and $a_{-1}=b$.  The endpoint conditions and
\eqref{eq:selected-blocks} imply, for every
$v=x_{i_1}^{\sigma_1}\cdots x_{i_t}^{\sigma_t}\in\mathcal V_m$,
\begin{equation}\label{eq:selected-multipliers}
 \norm{T(v(u^{(m)}))-a_{\sigma_1}v(u^{(m)})}_2
 =\norm{D(v(u^{(m)}))-a_{\sigma_1}}_2<2/m.
\end{equation}
The equality uses $T(v)=D(v)v$ and the $2$-norm isometry of right
multiplication by the unitary $v$.

\medskip\noindent
\emph{The exact family in the ultrapower.}
Fix a free ultrafilter $\omega_1$ on the tuple index $m$.  Set
$u_i^{(m)}=1$ when $m<i$, and define
\[
 N=M^{\omega_1},\qquad
 g_i=[(u_i^{(m)})_{m\ge1}]_{\omega_1}\quad(i\ge1),
 \qquad R=T^{\omega_1}.
\]
Lemma~\ref{lem:ultrapower} makes $R$ a bounded linear map on $N$,
with $\norm{R}\le C$, and supplies its bounded-ball $2$-continuity.
Each $g_i$ is unitary.  Every fixed nonidentity reduced word $q$ lies
in $\mathcal Q_m$ for all sufficiently large $m$, so
\eqref{eq:selected-moments} gives
\begin{equation}\label{eq:exact-haar-family}
 \tau^{\omega_1}(q(g_1,g_2,\ldots))=0.
\end{equation}
These are precisely the joint moments of a countable free Haar
family: powers of one generator are centered, and expansion of
alternating centered Laurent polynomials in different generators
reduces to nonidentity words.  The corresponding freeness of their
generated von Neumann algebras follows by bounded $2$-norm
approximation.

Identify $a,b\in M$ with their constant images in $N$.  For every
nonempty word $w=g_{i_1}^{\sigma_1}\cdots g_{i_t}^{\sigma_t}$ with
pairwise distinct unsigned indices,~\eqref{eq:selected-multipliers}
holds for all sufficiently large tuple indices.  Hence it becomes the
exact identity
\begin{equation}\label{eq:first-letter-multipliers}
 R(w)=
 \begin{cases}
  aw,&\sigma_1=1,\\
  bw,&\sigma_1=-1.
 \end{cases}
\end{equation}
These are all the hypotheses of Theorem~\ref{thm:rigidity}, applied
to $(N,\tau^{\omega_1})$ and $R$.  That theorem yields $a=b$ in $N$.
The constant embedding preserves the faithful trace, so this also
means $a=b$ in $M$, contrary to their choice.  Thus $Z_\infty$ is
deterministic, and the earlier argument proves
\eqref{eq:liouville-conclusion}.
\end{proof}

Only the scalar joint moments of the $g_i$ and the identities
\eqref{eq:first-letter-multipliers} enter the rigidity theorem.
Freeness of these unitaries from $a$ or $b$ is not asserted or needed.
